% MAS1004 과제 1: 짧은 보고서, 한국어 양식.
%
% 직접 쓴다. 어떤 단계에서도 언어 모델을 쓰지 않는다. 다 쓴 뒤 다듬을 때도,
% Overleaf 안의 AI 기능도 쓰지 않는다.
% 한 쪽. 한 쪽을 넘으면 끝에 빨간 글씨가 나온다.
% XeLaTeX로 컴파일한다 (Overleaf: Menu, Compiler, XeLaTeX).

\documentclass[11pt]{article}
\usepackage[onepage,ko]{mas1004}

\reporttitle{과제 1: 짧은 보고서}
\studentname{이름}
\studentid{학번}

\begin{document}
\makeheader

\question{1. 모델이 구분하는 것과 그것을 고른 이유}

여기에 쓴다.

\question{2. 모델이 저지르는 가장 나쁜 실수. 어떤 이미지였고, 무엇이라고
답했고, 무엇 때문에 그렇게 답했다고 생각하는지}

여기에 쓴다.

\question{3. 코딩 에이전트가 틀린 것 하나. 에이전트가 무엇이라고 했고, 실제로는
무엇이 맞았고, 무엇을 보고 확인하게 되었는지}

여기에 쓴다.

\question{4. 내일 처음부터 다시 한다면 다르게 할 것}

여기에 쓴다.

\end{document}
